% !PRINTED 236
\weberpage{236}

%<<CONT
identischen Substitution, so muss $\varphi^2 = 1$ sein, und dies ist auch ausreichend. Wir k\"onnen hier durch Transformation $\varphi$ auf die Form $\begin{pmatrix} \alpha, & 0 \\ 0, & -\alpha^{-1} \end{pmatrix}$ bringen und erhalten als Bedingung $\alpha^2 = \pm 1$, und demnach zwei Formen der Substitution $\varphi$:
\[
\varphi' = \begin{pmatrix} i, & 0 \\ 0, & i \end{pmatrix}, \quad
\varphi'' = \begin{pmatrix} 1, & 0 \\ 0, & -1 \end{pmatrix},
\]
von denen der erste die \gesperrt{Inversion}, der zweite die \gesperrt{Spiegelung} genannt wird. Uebertr\"agt man sie nach \S.\,50, (1) auf eine rechtwinkelige Coordinatentransformation, so bedeutet $\varphi'$ die gleichzeitige Vorzeichen\"anderung aller drei Coordinaten, $\varphi''$ die Vertauschung von $x$ mit $-x$. Es ergeben sich hieraus die beiden Gruppen der zweiten Art:
\weq{2}{C_1' = \begin{pmatrix} i^h, & 0 \\ 0, & i^h \end{pmatrix}, \quad C_1'' = \begin{pmatrix} 1, & 0 \\ 0, & (-1)^h \end{pmatrix}, \quad h = 0, 1.}

II. Es sei $P$ eine cyklische Gruppe $C_n$, $n > 1$.
Setzen wir
\weq{3}{c = \begin{pmatrix} \varepsilon, & 0 \\ 0, & \varepsilon^{-1} \end{pmatrix}, \quad \varepsilon = e^{\frac{\pi i}{n}},}
so ist nach \S.\,55:
\weq{4}{C_n = 1, c, c^2, \ldots, c^{n-1},}
worin
\[
c^s = \begin{pmatrix} \varepsilon^s, & 0 \\ 0, & \varepsilon^{-s} \end{pmatrix};
\]
wenn nun
\weq{5}{\varphi = \begin{pmatrix} \alpha, & \beta \\ \gamma, & \delta \end{pmatrix}, \quad \alpha\delta - \beta\gamma = -1}
ist, so bilden wir zun\"achst
\weq{6}{\varphi^{-1} c \varphi = \begin{pmatrix} \alpha\delta\varepsilon - \beta\gamma\varepsilon^{-1}, & \beta\delta(\varepsilon - \varepsilon^{-1}) \\ -\alpha\gamma(\varepsilon - \varepsilon^{-1}), & \alpha\delta\varepsilon^{-1} - \beta\gamma\varepsilon \end{pmatrix},}
\weq{7}{\varphi^2 = \begin{pmatrix} \alpha^2 + \beta\gamma, & (\alpha + \delta)\beta \\ (\alpha + \delta)\gamma, & \beta\gamma + \delta^2 \end{pmatrix}.}
Diese beiden Substitutionen m\"ussen in der Reihe (4) vorkommen, und es muss also $\beta\delta = 0$, $\alpha\gamma = 0$ sein; also m\"ussen entweder $\beta$ und $\gamma$ oder $\alpha$ und $\delta = 0$ sein. Ist $\beta = \gamma = 0$, so folgt aus (7), dass $\alpha$ von der Form $\varepsilon^{1/2\,r}$, also
\[
\varphi = \begin{pmatrix} e^{\frac{\pi i}{2n}\,r}, & 0 \\ 0, & -e^{-\frac{\pi i}{2n}\,r} \end{pmatrix}
\]
%CONT>>
